\section{Saved results on the required instances}
Benchmark gaps use published BKS values. For the team instance the reference is the best value
already found in the project, $20{,}724.21095$, rather than a proven optimum.
Mean gap is $100(\text{mean cost}-\text{reference})/\text{reference}$.

\begin{table}[h]\centering\small\setlength{\tabcolsep}{4pt}
\begin{tabular}{@{}llrrrrr@{}}
\toprule
Instance & Alg. & Best & Mean & SD & Time (s) & Mean gap (\%) \\
\midrule
A-n32-k5 & CW & 842.0 & -- & -- & $<0.01$ & 7.40 \\
A-n32-k5 & GA & 827.0 & 827.0 & 0.00 & 30.00 & 5.48 \\
A-n32-k5 & TS & 784.0 & 784.0 & 0.00 & 30.00 & 0.00 \\
A-n45-k6 & CW & 997.0 & -- & -- & $<0.01$ & infeasible \\
A-n45-k6 & GA & 944.0 & 954.9 & 11.51 & 30.00 & 1.15 \\
A-n45-k6 & TS & 944.0 & 944.0 & 0.00 & 30.00 & 0.00 \\
A-n60-k9 & CW & 1,408.0 & -- & -- & $<0.01$ & 3.99 \\
A-n60-k9 & GA & 1,382.0 & 1,396.2 & 6.94 & 60.00 & 3.12 \\
A-n60-k9 & TS & 1,354.0 & 1,354.0 & 0.00 & 60.00 & 0.00 \\
A-n80-k10 & CW & 1,840.0 & -- & -- & $<0.01$ & 4.37 \\
A-n80-k10 & GA & 1,818.0 & 1,818.0 & 0.00 & 60.00 & 3.12 \\
A-n80-k10 & TS & 1,763.0 & 1,763.6 & 0.97 & 60.00 & 0.03 \\
Team (7273) & CW & 22,573.3 & -- & -- & $<0.01$ & infeasible \\
Team (7273) & GA & 21,166.7 & 21,352.2 & 114.24 & 60.00 & 3.03 \\
Team (7273) & TS & 20,724.2 & 20,725.5 & 4.05 & 60.00 & 0.01 \\
\bottomrule
\end{tabular}

\caption{GA and TS each have ten runs. SD is sample standard deviation. Time is mean total
runtime, not time to best. CW runs once, so mean and SD are omitted. Fleet-infeasible CW costs
are diagnostic and receive no gap.}
\end{table}

TS reaches the benchmark reference in all ten runs on A-n32-k5, A-n45-k6 and A-n60-k9. On
A-n80-k10 it reaches 1763 in six runs, with mean 1763.6. GA has a higher mean on every benchmark.
On A-n45-k6, GA sometimes reaches the reference, but less consistently.

On the team instance TS averages about 20,725.5, compared with 21,352.2 for GA. CW needs 19
routes when only 18 are allowed. Each route fits capacity, but the plan is still invalid.
These observations concern the implemented methods and settings, not every possible GA or TS.

\subsection*{Tiny instances: exact answers}
\begin{table}[h]\centering\small
\begin{tabular}{@{}lrrcc@{}}
\toprule
Instance & $n$ & Exact cost & GA hits & TS hits \\
\midrule
tiny-n7-0 & 7 & 3180.33 & 10/10 & 10/10 \\
tiny-n7-1 & 7 & 2753.56 & 10/10 & 10/10 \\
tiny-n7-2 & 7 & 2873.84 & 10/10 & 10/10 \\
tiny-n8-0 & 8 & 3069.92 & 10/10 & 10/10 \\
tiny-n8-1 & 8 & 3535.17 & 10/10 & 10/10 \\
tiny-n8-2 & 8 & 2987.33 & 10/10 & 10/10 \\
tiny-n9-0 & 9 & 4212.50 & 10/10 & 10/10 \\
tiny-n9-1 & 9 & 4359.07 & 10/10 & 10/10 \\
tiny-n9-2 & 9 & 5307.12 & 10/10 & 10/10 \\
\bottomrule
\end{tabular}

\caption{Stored exact costs and one-second runs reaching each optimum.
Cached dynamic-programming and MILP exact values agree.}
\end{table}

Each method reaches the exact optimum in all 90 tiny runs. This supports correctness on small
cases; it does not guarantee optimality on larger ones. The independent Python checker verifies
coverage, loads, fleet size and recomputed cost without using the C++ evaluator.
